\section{Complementary Reductions of One Transducer Theory}\label{sec:reductions}
The quantum reading retains preparations, internal operations, phase, and registered measure. The geometric reading retains causal accessibility, accumulated transport, and clock comparisons. They are proposed reductions of one admitted process, not exhaustive descriptions of quantum theory and relativity. A valid pair of reductions must declare which distinctions each retains and discards.

\section{The Common-Invariant Criterion}\label{sec:commoninvariant}
For a specified transducer $\mathfrak T$, seek a native invariant $\mathcal J(\mathfrak T)$ and maps $R_Q,R_G$ deriving quantum and geometric observables from the same input and admission law. A common graph is not enough if the two models require independent adjustable structures. A common invariant is a substantive bridge; full compatibility also requires dynamics, observable composition, and controlled approximation in overlapping regimes.

The receipt algebra, relative orientation, phase holonomy, and registered winding are candidate inputs. A native relational-complexity functional would supply another candidate only after its domain, invariance, and compatibility with these readouts have been established. Mode supplies a precise presentation character, while the transport and temporal maps must show which of these structures they actually preserve.


One intermediate target is explicit: a common state representation, invariant form, admitted action family, and registered readouts, written $(\mathcal X,\mathsf G,\mathfrak A,\mathfrak O)$. For a fixed-form linear conservative sector its generators obey (CD2). The generator, observables, and any field-address locality must be supplied together; a form that can be manufactured by finite-group averaging is not by itself a common physical foundation. Sections~\ref{sec:conservativedynamics}--\ref{sec:nativedynamics} develop the comparison and its failure tests.

\paragraph{A common domain need not identify its endpoint objects.}
The incidence span (IR4) offers a more general target than a direct equality of two registers. History and analyzer presentations may remain different while participating in one native incidence. Its projected maps and lifted actions must commute with the actual admitted transport. Proposition~\ref{prop:reconciliationcriterion} then tests whether their orientation comparisons are compatible. Selecting references makes the comparison expressible; it does not supply any of those maps.

\section{Event Orientation and Transport Structure}\label{sec:orientation}
For local event operators $E_s,E_t$ and a compatible admitted carrier map $C_{ts}$, the original test is
% Original equation number(s) 47; expression unchanged.
\setcounter{equation}{46}
\begin{equation}
 C_{ts}E_s=\sigma_{ts}E_tC_{ts},
 \qquad\sigma_{ts}\in\{+1,-1\}.
\label{eq:eventorientation}
\end{equation}

Composite transports multiply their signs. A local convention $E_v'=\eta_vE_v$ changes an edge sign to $\sigma'_{ts}=\eta_s\eta_t\sigma_{ts}$, leaving closed-loop products invariant.

Where an explicit map identifies the chosen event-orientation representation with the registered modal character, $\sigma_{ts}$ can express modal preservation or reversal. Until that map is supplied, it remains the sign of the stated intertwining equation. Mode, deck residue, Born grading, and $r_1$ winding retain their distinct domains. This is the concrete interface through which event-preserving transport can be tested rather than inferred from a matching sign.

\paragraph{Naturality at the declared interfaces.}
The Program~1 receipts require a state-dependent $C_4$ coefficient for the neutral face axis: the common exchange has different reported signs on its two fixed faces. The correct order-ten face-realization kernels nevertheless transport naturally, as confirmed directly by the native incidence replay. Program~2 uses the normal within-sector subgroup of the separate analyzer-sector extension. These results close their respective finite vertical-transport questions; they do not identify the three programs' orientation objects \citep{cave2026p4receipts,cave2026draft6replays}.

For Program~3, neither the literal phase-groupoid automorphism nor the split signed-axis group supplies the missing common incidence and action on the correlation. Its VN03/VN04 interface remains open. The appropriate target is native compatibility of actions, possibly only up to a common sign convention, rather than an absolute positive history label. Appendix~\ref{app:referenceaudit} records the exact audit boundaries.

\section{Action, Phase, and Planck's Constant}\label{sec:hbar}
Primitive informative action is a realized distinction; mechanical action $S$ is a dimensionful history functional. An admitted scalar phase representation can obey
% Original equation number(s) 48; expression unchanged.
\setcounter{equation}{47}
\begin{equation}
\chi(\gamma_2\circ\gamma_1)=\chi(\gamma_2)\chi(\gamma_1),
\qquad\chi(\gamma)\in U(1).
\end{equation}

Then $\chi=e^{i\theta}$ gives an additive phase modulo $2\pi$. A real-valued lift needs branch or covering information. The finite three-phase coordinate of Section~\ref{sec:mode} is not automatically this $U(1)$ phase.

A useful intermediate target is an additive native functional $S_{\rm nat}$ and a structure-selected relation $\chi(\gamma)=\exp(i\kappa S_{\rm nat}(\gamma))$. A physical energy-frequency relation then additionally needs calibrated dynamics and time. The conventional $U(t)=e^{-iHt/\hbar}$ relation supplies the comparison target, not a graph-derived value of $\hbar$ \citep{bipm}. No action or clock scale is changed by the definition of Mode.
